\section*{Asesmen}
\addcontentsline{toc}{section}{Asesmen}

Bab XVII tidak memiliki bobot tersendiri. Asesmennya terintegrasi
dalam UAS (16\%) dan Project Final (20\%). Fokusnya adalah
kemampuan menghubungkan konsep, menganalisis, dan merefleksikan.

\begin{table}[H]
\centering
\caption{Instrumen asesmen integratif}
\label{tab:ch17-asesmen}
\begin{tabularx}{\textwidth}{@{}lXl@{}}
\toprule
\textbf{Instrumen} & \textbf{Indikator} & \textbf{Bobot} \\
\midrule
UAS & Integrasi konsep dan studi kasus & 16\% \\
Project Final & Pipeline lengkap, demo, laporan & 20\% \\
Portofolio & Bukti belajar dan refleksi & bagian proyek \\
\bottomrule
\end{tabularx}
\end{table}

Kriteria kelulusan: mahasiswa dapat menjelaskan seluruh pipeline,
menganalisis kesalahan, mempertahankan keputusan desain, dan
menyusun rencana pengembangan diri.
